
\begin{obeactivity}
\textbf{Aktivitas 10.1: Simulasi RSA dengan Bilangan Prima Kecil}
1. Pilih dua bilangan prima kecil, misal $p=3$ dan $q=11$.
2. Hitung $n$ dan $\phi(n)$.
3. Pilih $e$ yang relatif prima terhadap $\phi(n)$.
4. Hitung $d$ menggunakan metode coba-coba atau Algoritma Euclidean.
5. Cobalah mengenkripsi pesan $m=5$ dan dekripsi kembali hasilnya.
6. Diskusikan apa yang terjadi jika $e$ tidak relatif prima terhadap $\phi(n)$.
\end{obeactivity}

Langkah-langkah pada Aktivitas 10.1 dapat dijalankan langsung melalui program
pada Listing~\ref{lst:rsa}. Program itu membangkitkan kunci, mengenkripsi, dan
mendekripsi pesan untuk parameter Aktivitas 10.1 maupun parameter pada contoh
terhitung Bab 10, sehingga setiap angka dapat diperiksa satu per satu.

\lstinputlisting[
    language=Python,
    caption={RSA: pembangkitan kunci, enkripsi, dan dekripsi},
    label={lst:rsa}
]{\codefile{python/rsa.py}}

Program yang sama tersedia dalam C++ (\texttt{code/cpp/rsa.cpp}) dan MATLAB
(\texttt{code/matlab/rsa.m}).
